\documentclass[reqno,11pt]{amsart}

\usepackage[margin=1in]{geometry}
\usepackage[T1]{fontenc}
\usepackage{lmodern}
\usepackage{microtype}
\usepackage{amsmath,amssymb,amsthm,mathtools}
\usepackage[dvipsnames]{xcolor}
\usepackage[most]{tcolorbox}
\usepackage{enumitem}
\usepackage[hidelinks]{hyperref}
\hypersetup{pdftitle={Lecture 7, MATH 327: Cosets and Partitions},pdfauthor={Manish Patnaik},pdfsubject={MATH 327 course notes}}

\definecolor{courseblue}{HTML}{1F4E79}
\definecolor{lightblue}{HTML}{EAF2F8}
\definecolor{darkgray}{HTML}{333333}
\definecolor{definitiongreen}{HTML}{EEF8EE}
\definecolor{definitionborder}{HTML}{8FB58F}
\definecolor{lightpink}{HTML}{FCE8EF}
\definecolor{warningyellow}{HTML}{FFF7D6}
\definecolor{warningborder}{HTML}{D5A500}
\setlength{\parindent}{0pt}
\setlength{\parskip}{0.65em}
\setlist[itemize]{topsep=0.25em,itemsep=0.15em}

\makeatletter
\def\@tocline#1#2#3#4#5#6#7{\relax
  \ifnum #1>\c@tocdepth\else
    \par\addpenalty\@secpenalty\addvspace{#2}\begingroup\hyphenpenalty\@M
    \@ifempty{#4}{\@tempdima\csname r@tocindent\number#1\endcsname\relax}{\@tempdima#4\relax}
    \parindent\z@\leftskip#3\relax\advance\leftskip\@tempdima\relax
    \rightskip\@pnumwidth plus4em\parfillskip-\@pnumwidth
    #5\leavevmode\hskip-\@tempdima\ifcase #1\or\or\hskip 1em\or\hskip 2em\else\hskip 3em\fi
    #6\nobreak\relax\dotfill\hbox to\@pnumwidth{\@tocpagenum{#7}}\par\nobreak\endgroup
  \fi}
\makeatother

\theoremstyle{definition}
\newtheorem{definition}{Definition}[section]
\tcolorboxenvironment{definition}{enhanced,breakable,colback=definitiongreen,colframe=definitionborder,boxrule=0.5pt,arc=1mm,left=2mm,right=2mm,top=1mm,bottom=1mm,before skip=0.8em,after skip=0.8em}
\newtheorem{example}{Example}[section]
\theoremstyle{plain}
\newtheorem{theorem}{Theorem}[section]
\newtheorem{lemma}{Lemma}[section]
\newtheorem{proposition}{Proposition}[section]
\theoremstyle{remark}
\newtheorem*{remark}{Remark}
\newtheorem{exercise}[theorem]{Exercise}
\newtcolorbox{warningbox}{enhanced,breakable,colback=warningyellow,colframe=warningborder,coltitle=darkgray,fonttitle=\bfseries,title={Warning},boxrule=0.7pt,arc=2pt,left=6pt,right=6pt,top=5pt,bottom=5pt,before skip=0.8em,after skip=0.8em}
\tcolorboxenvironment{exercise}{enhanced,colback=lightpink,colframe=WildStrawberry,boxrule=0.7pt,arc=2pt,left=6pt,right=6pt,top=5pt,bottom=5pt}
\numberwithin{equation}{section}
\newcommand{\Z}{\mathbb{Z}}

\begin{document}
\begin{center}
  \colorbox{lightblue}{\begin{minipage}{0.94\textwidth}
    \vspace{0.55em}
    \begin{tabular*}{\textwidth}{@{\extracolsep{\fill}}lr}
      {\color{courseblue}\bfseries\LARGE Lecture 7} & {\color{courseblue}\bfseries\LARGE MATH 327}\\[0.35em]
      & \textbf{Date: Sep. 18, 2026}\\
      & \textbf{Instructor: Manish Patnaik}
    \end{tabular*}
    \vspace{0.45em}
  \end{minipage}}
\end{center}
\setcounter{tocdepth}{1}
\tableofcontents

\section{Normality revisited}

Recall that $N\leq G$ is normal if $gNg^{-1}\subseteq N$ for every $g\in G$.

\begin{proposition}
For a subgroup $N\leq G$, the following are equivalent:
\begin{enumerate}[label=\textbf{(\arabic*)}]
  \item $gNg^{-1}\subseteq N$ for every $g\in G$;
  \item $gNg^{-1}=N$ for every $g\in G$;
  \item $gN=Ng$ for every $g\in G$.
\end{enumerate}
\end{proposition}

\begin{proof}
Suppose (1) holds. Applying it to $g^{-1}$ gives
$g^{-1}Ng\subseteq N$. If $n\in N$, then $g^{-1}ng\in N$, and hence
$n=g(g^{-1}ng)g^{-1}\in gNg^{-1}$. Thus $N\subseteq gNg^{-1}$ and (2)
follows. Multiplying the equality $gNg^{-1}=N$ on the right by $g$ gives
$gN=Ng$, proving (3). Conversely, multiplying $gN=Ng$ on the right by
$g^{-1}$ gives $gNg^{-1}=N$.
\end{proof}

For example, in
\[
  S_3=\{1,x,x^2,y,xy,yx\},
  \qquad x^3=y^2=1,
  \qquad yx=x^2y,
\]
the subgroup $N=\{1,x,x^2\}$ satisfies $yNy^{-1}=N$, and hence is normal.
This calculation will reappear naturally in terms of cosets.

Writing out the calculation used in the previous lecture,
\begin{align*}
  yNy^{-1}
  &=(yN)y^{-1}\\
  &=\{y,yx,yx^2\}y^{-1}\\
  &=\{1,yxy,yx^2y\}\\
  &=\{1,x^2,x\}=N.
\end{align*}
Multiplying the equality $yNy^{-1}=N$ on the right by $y$ also gives
$yN=Ny$.

\section{Left and right cosets}

\begin{definition}
Let $H\leq G$ and $g\in G$. The \emph{left coset} and \emph{right coset}
of $H$ represented by $g$ are
\begin{equation}
  gH=\{gh:h\in H\},
  \qquad
  Hg=\{hg:h\in H\}.
\end{equation}
The element $g$ is called a \emph{coset representative}.
\end{definition}

If $g\in H$, then $gH=H=Hg$. A representative is not unique: every element
of a coset represents that same coset.

\begin{example}[Cosets in $S_3$]
Let $H=\{1,x,x^2\}\leq S_3$. Then
\[
  1H=xH=x^2H=H,
\]
whereas
\[
  yH=xyH=yxH=\{y,xy,yx\}.
\]
Thus the set of left cosets has two elements:
\[
  S_3/H=\{H,yH\}.
\]
Since $H$ is normal, these are also the two right cosets.
\end{example}

The individual computations are
\begin{align*}
  1H&=\{1,x,x^2\},\\
  xH&=\{x,x^2,x^3\}=\{1,x,x^2\}=H,\\
  x^2H&=\{x^2,x^3,x^4\}=\{1,x,x^2\}=H,\\
  yH&=\{y,yx,yx^2\}=\{y,yx,xy\},\\
  xyH&=\{xy,xyx,xyx^2\}=\{xy,y,yx\},\\
  yxH&=\{yx,yx^2,yx^3\}=\{yx,xy,y\}.
\end{align*}
Thus every element of $S_3$ lies in exactly one of the two cosets $H$ and
$yH$.

\begin{lemma}[Equality of cosets]
Let $H\leq G$ and $g_1,g_2\in G$. Then
\begin{equation}
  g_1H\cap g_2H\neq\varnothing
  \quad\Longleftrightarrow\quad
  g_1H=g_2H
  \quad\Longleftrightarrow\quad
  g_2^{-1}g_1\in H.
\end{equation}
\end{lemma}

\begin{proof}
First consider the special case $g_2=e$. We have
\[
  gH\cap H\neq\varnothing
  \quad\Longleftrightarrow\quad
  g\in H
  \quad\Longleftrightarrow\quad
  gH=H.
\]
Indeed, if $g\in H$, closure under multiplication gives $gH\subseteq H$;
and for $h\in H$, the expression $h=g(g^{-1}h)$ shows $H\subseteq gH$.
Conversely, if $x\in gH\cap H$, write $x=gh$ with $h\in H$. Then
$g=xh^{-1}\in H$.

For the general case, left multiplication by $g_2^{-1}$ reduces it to the
special case:
\[
  g_1H\cap g_2H\neq\varnothing
  \quad\Longleftrightarrow\quad
  g_2^{-1}g_1H\cap H\neq\varnothing
  \quad\Longleftrightarrow\quad
  g_2^{-1}g_1\in H.
\]
The same reduction shows that this is equivalent to $g_1H=g_2H$.
\end{proof}

For completeness, the reduction may also be written elementwise.  If
$x\in g_1H\cap g_2H$, then $x=g_1h_1=g_2h_2$ for some $h_1,h_2\in H$.
Therefore
\[
  g_2^{-1}g_1=h_2h_1^{-1}\in H.
\]
Conversely, if $g_2^{-1}g_1=h\in H$, then $g_1=g_2h\in g_2H$, so the two
cosets intersect and hence are equal.

\begin{proposition}
The collection of left cosets $\{gH:g\in G\}$ is a partition of $G$.
That is, every element of $G$ belongs to a left coset, and two left cosets
are either equal or disjoint.
\end{proposition}

\begin{proof}
Every $g\in G$ lies in $gH$ because $g=ge$ and $e\in H$. The preceding
lemma shows that any two cosets that intersect are equal.
\end{proof}

\begin{definition}[Coset space]
The set of left cosets of $H$ in $G$ is denoted by
\begin{equation}
  G/H=\{gH:g\in G\}.
\end{equation}
\end{definition}

\section{Example: cosets of \texorpdfstring{$8\Z$}{8Z}}

Consider the additive group $G=\Z$ and its subgroup $H=8\Z$. For
$n\in\Z$, the coset represented by $n$ is
\begin{equation}
  n+8\Z=\{n+8k:k\in\Z\}.
\end{equation}
It consists precisely of the integers congruent to $n$ modulo $8$.

For example,
\begin{align*}
  0+8\Z&=\{\ldots,-16,-8,0,8,16,24,\ldots\},\\
  1+8\Z&=\{\ldots,-15,-7,1,9,17,25,\ldots\},\\
  2+8\Z&=\{\ldots,-14,-6,2,10,18,26,\ldots\},\\
  7+8\Z&=\{\ldots,-9,-1,7,15,23,31,\ldots\}.
\end{align*}
The remaining four cosets are obtained in the same way.  Arranging integers
by their remainder on division by $8$ gives eight disjoint columns, one for
each coset.

There are eight distinct cosets:
\begin{equation}
  \Z/8\Z
  =\{0+8\Z,1+8\Z,\dots,7+8\Z\}.
\end{equation}
For instance,
\[
  0+8\Z=8+8\Z,
  \qquad
  1+8\Z=17+8\Z.
\]
Writing $\overline n=n+8\Z$, we may abbreviate the coset space as
\begin{equation}
  \Z/8\Z=\{\overline0,\overline1,\dots,\overline7\}.
\end{equation}

More generally, the equality-of-cosets lemma becomes the familiar congruence
criterion
\begin{equation}
  a+8\Z=b+8\Z
  \quad\Longleftrightarrow\quad
  a-b\in8\Z
  \quad\Longleftrightarrow\quad
  a\equiv b\pmod 8.
\end{equation}

\begin{remark}
At this stage $G/H$ denotes a set of cosets. When $H\trianglelefteq G$, these
cosets will inherit a well-defined group operation; this leads to quotient
groups.
\end{remark}

\end{document}
